% !TEX program = pdflatex
% LNCS/ICONIP-style summary paper. Layout follows the referenced main.tex.
\documentclass[runningheads]{llncs}

\usepackage[T1]{fontenc}
\usepackage{graphicx}
\usepackage{booktabs}
\usepackage{multirow}
\usepackage{array}
\usepackage{amsmath}
\usepackage{amssymb}
\usepackage{url}

\urlstyle{rm}

\title{Wavelet-Supervised Test-Time Prototype Adaptation for Zero-Shot Anomaly Detection}
\titlerunning{Wavelet-Supervised Prototype Adaptation for ZSAD}
\author{Anonymous Author(s)}
\authorrunning{Anonymous Author(s)}
\institute{Anonymous Institution}

\newcommand{\simi}{\operatorname{sim}}
\newcommand{\normop}{\operatorname{norm}}

\begin{document}
\maketitle

\begin{abstract}
Zero-shot industrial anomaly detection requires image-level anomaly recognition and pixel-level anomaly localization without target-category training images or anomaly annotations. Existing CLIP-based methods usually rely on fixed normal/abnormal text prototypes, which cannot fully capture instance-specific defect appearances and may confuse true defects with normal material textures or structural boundaries. We propose Wavelet-Supervised Test-Time Prototype Adaptation (WPTA), which uses Haar wavelet cues as patch evidence reliability for test-time prototype adaptation rather than as an additive anomaly map. WPTA freezes CLIP, constructs boundary-aware wavelet reliability on the CLIP patch feature grid, combines it with the initial semantic anomaly prior to select visual normal/abnormal anchors, and calibrates prototypes through confidence-gated conservative updates. Controlled ablations on MVTec and VisA show that direct wavelet fusion degrades localization, while semantic-spectral evidence selection and conservative calibration improve test-time prototype adaptation. Across five industrial benchmarks, the calibrated inference system improves the fixed AnomalyCLIP baseline by $+1.0/+4.5/+3.4/+2.8$ on average in P-AUROC/P-AUPRO/I-AUROC/I-AP.
\keywords{Zero-shot anomaly detection \and CLIP \and Test-time adaptation \and Wavelet reliability}
\end{abstract}

\section{Introduction}

Zero-shot anomaly detection (ZSAD) aims to detect and localize anomalies on unseen target categories without using target-domain training images. CLIP-based ZSAD methods have made this setting practical by comparing visual features with normal and abnormal text prototypes. However, a fixed abnormal prototype only provides a generic anomaly direction. It does not necessarily match the concrete defect appearance in the current image, such as subtle scratches, missing edges, surface contamination, or texture damage. As a result, genuine defects may be underestimated, while normal boundaries, reflections, or material textures may be activated as false positives.

Frequency and wavelet cues provide useful but risky complementary information. Many industrial defects manifest as local high-frequency disturbances, and Haar wavelets can separate high-frequency details from low-frequency structures on feature grids. Nevertheless, high-frequency response is not equivalent to anomaly evidence. Normal object contours, reflective surfaces, and fine-grained textures can also produce strong wavelet responses. In controlled tests, direct wavelet fusion lowers pixel-level (AUROC, AUPRO) on MVTec from $(91.2,83.2)$ to $(88.7,80.4)$, and on VisA from $(95.5,86.7)$ to $(94.6,85.1)$.

We therefore reposition wavelet cues from final-map scoring to evidence reliability for test-time prototype adaptation. WPTA first computes an initial semantic anomaly prior $S_0$ using fixed normal/abnormal prototypes. It then applies Haar decomposition to the CLIP patch feature grid, estimates high-frequency texture energy and low-frequency structural edges, and constructs a boundary-aware reliability map. The semantic prior $S_0$ provides the anomaly direction, while the wavelet reliability $W$ constrains whether local evidence is trustworthy. The selected visual anchors are used to conservatively calibrate prototypes without backpropagation or CLIP parameter updates.

Our evidence is organized at two levels. First, controlled MVTec/VisA ablations validate the WPTA mechanism: semantic-only prototype adaptation improves over fixed prototypes, wavelet-guided reliability further improves localization, and conservative update gives additional gains. Second, five industrial benchmarks report the system-level effect of the calibrated inference stack. Since the five-dataset final systems use dataset-specific calibration modules, we explicitly separate system-level gains from the causal evidence for WPTA.

The contributions are summarized as follows:
\begin{itemize}
  \item We propose WPTA, which turns Haar wavelet cues into patch evidence reliability for test-time prototype adaptation rather than direct final-map fusion.
  \item We design semantic-spectral evidence selection, jointly using the semantic anomaly prior and boundary-aware wavelet reliability to construct visual normal/abnormal anchors.
  \item We perform conservative, training-free prototype calibration with frozen CLIP and validate the mechanism through controlled MVTec/VisA ablations.
  \item We report five industrial benchmark results with a configuration audit that separates calibrated-system gains from WPTA mechanism evidence.
\end{itemize}

\begin{figure}[t]
  \centering
  \includegraphics[width=0.92\textwidth]{figures/figure1_motivated_example_mvtec_cable_vector_layout.pdf}
  \caption{Motivating example. Wavelet cues are useful as evidence reliability for prototype adaptation rather than as direct anomaly scores.}
  \label{fig:motivation}
\end{figure}

\section{Related Work}

\paragraph{CLIP-based zero-shot anomaly detection.}\mbox{}\\
\noindent
CLIP~\cite{radford2021clip} provides a general image--text alignment foundation for open vocabulary recognition. WinCLIP~\cite{jeong2023winclip} performs window-level CLIP matching for zero-/few-shot anomaly classification and segmentation. AnomalyCLIP~\cite{zhou2024anomalyclip} learns object-agnostic normal/abnormal prompts, while AdaCLIP~\cite{cao2024adaclip} adapts CLIP with hybrid learnable prompts. CLIP-AD~\cite{chen2024clipad} improves zero-shot anomaly detection with language-guided staged dual-path modeling. These methods show the value of CLIP semantic prototypes, but most of them use fixed or trained prototypes without training-free calibration from local evidence in the current test image.

\paragraph{Test-time adaptation.}\mbox{}\\
\noindent
Test-time adaptation uses information from the test sample to mitigate distribution shifts. TPT~\cite{shu2022tpt} performs test-time prompt tuning for vision-language models through confidence selection and view consistency. Anomaly detection is more challenging because abnormal regions are sparse, unknown, and often adjacent to normal boundaries or textures. WPTA avoids updating CLIP parameters or adapting from the whole image. It aggregates only patch evidence jointly supported by the semantic prior and wavelet reliability.

\paragraph{Frequency and wavelet cues.}\mbox{}\\
\noindent
Wavelet representations decompose low-frequency structures and high-frequency details at multiple resolutions~\cite{daubechies1988orthonormal,mallat1989theory}. In industrial anomaly localization, high-frequency disturbances may indicate scratches, cracks, or texture damage. However, direct frequency scoring is vulnerable to normal boundaries and textures. WPTA limits wavelet cues to evidence reliability and leaves final anomaly decisions to calibrated CLIP prototypes.

\section{Method}

\subsection{Problem Setting and Overall Framework}

Given a test image $x$, ZSAD predicts a pixel-level anomaly map $M$ and an image-level anomaly score $s_{\mathrm{img}}$. The frozen CLIP image encoder extracts patch features $\{f_i\}_{i=1}^{N}$, and the text encoder provides a normal prototype $t_n$ and an abnormal prototype $t_a$. WPTA performs inference without target-domain training images, anomaly annotations, backpropagation, or CLIP parameter updates.

The method has four steps, as shown in Fig.~\ref{fig:architecture}. First, initial semantic scoring computes a patch-level anomaly prior $S_0$. Second, boundary-aware Haar wavelet reliability estimates a reliability map $W$ on the CLIP patch feature grid. Third, semantic-spectral evidence selection aggregates visual normal/abnormal anchors. Fourth, conservative prototype calibration recomputes anomaly probabilities using calibrated prototypes.

\begin{figure}[t]
  \centering
  \includegraphics[width=0.92\textwidth]{figures/figure2_wpta_method_overview.pdf}
  \caption{Overall architecture of WPTA. Wavelet cues supervise evidence reliability for test-time prototype adaptation rather than directly producing the final anomaly map.}
  \label{fig:architecture}
\end{figure}

\subsection{Initial Semantic Anomaly Score}

For each patch feature $f_i$, WPTA computes the initial semantic anomaly prior by a softmax over similarities to normal and abnormal prototypes:
\begin{equation}
S_0(i)=
\frac{\exp(\simi(f_i,t_a)/\tau)}
{\exp(\simi(f_i,t_a)/\tau)+\exp(\simi(f_i,t_n)/\tau)} .
\end{equation}
Here $\simi(\cdot,\cdot)$ denotes normalized similarity and $\tau$ is the temperature. $S_0$ provides the semantic direction but is not treated as the final calibrated score.

\subsection{Boundary-Aware Haar Wavelet Reliability}

WPTA reshapes patch features into $F\in\mathbb{R}^{H\times W\times C}$ and applies one-level Haar DWT to obtain $LL$, $LH$, $HL$, and $HH$ components. The high-frequency texture energy is
\begin{equation}
HF(i)=\frac{1}{C}\sum_{c=1}^{C}\left(|LH_i^c|+|HL_i^c|+|HH_i^c|\right).
\end{equation}
Since high-frequency energy can also arise from normal structural boundaries, WPTA estimates a low-frequency structural edge response $LF_{\mathrm{edge}}$ from $LL$ and constructs
\begin{equation}
W(i)=\normop(HF(i))\cdot\left(1-\normop(LF_{\mathrm{edge}}(i))\right).
\end{equation}
Thus, high-frequency texture disturbance increases evidence reliability, while low-frequency structural edges suppress pseudo evidence.

\subsection{Semantic-Spectral Evidence Selection}

Abnormal evidence should be semantically close to the abnormal prototype and supported by reliable local disturbance. Normal evidence should be close to the normal prototype and not contaminated by abnormal textures or structure-boundary responses. We write the conceptual evidence weights as
\begin{equation}
q_a(i)=S_0(i)^\gamma\cdot\rho(W(i)),\qquad
q_n(i)=(1-S_0(i))^\gamma\cdot\rho(1-W(i)).
\end{equation}
Implementation uses top-$k$ selection, confidence thresholds, normalization, and confidence gates to control evidence quality.

\subsection{Visual Anchors and Conservative Calibration}

Selected evidence patches are aggregated into visual anchors:
\begin{equation}
v_a=\frac{\sum_i q_a(i)f_i}{\sum_i q_a(i)},\qquad
v_n=\frac{\sum_i q_n(i)f_i}{\sum_i q_n(i)}.
\end{equation}
The prototypes are then conservatively calibrated:
\begin{equation}
t'_a=\operatorname{Norm}\left((1-\alpha)t_a+\alpha v_a\right),\qquad
t'_n=\operatorname{Norm}\left((1-\beta)t_n+\beta v_n\right).
\end{equation}
The update strengths are controlled by confidence gates. In the default controlled setting, $\alpha_0=0.0$ and $\beta_0=0.01$, so abnormal evidence mainly gates the update, while the actual prototype displacement mainly affects the normal prototype. The final anomaly probability is recomputed by replacing $t_a,t_n$ with $t'_a,t'_n$ in the semantic scoring equation.

\section{Experiments}

\subsection{Experimental Setup}

\subsubsection{Datasets and Metrics}
\mbox{}\\
\noindent
We evaluate on five industrial anomaly detection benchmarks: MVTec AD~\cite{bergmann2019mvtec,bergmann2021mvtec}, VisA~\cite{zou2022visa}, MPDD~\cite{jezek2021deep}, BTAD~\cite{mishra2021vt}, and DTD-Synthetic~\cite{aota2023zero}. MVTec and VisA are used for controlled WPTA mechanism ablations. The five-dataset table reports system-level results of the calibrated inference stack.

We report Pixel-AUROC, Pixel-AUPRO, Image-AUROC, and Image-AP. Pixel-AUPRO is emphasized because it better reflects region-level anomaly localization quality. All results are deterministic reports without multi-seed confidence intervals or significance tests.

\subsubsection{Implementation Details}
All experiments use the AnomalyCLIP checkpoint \texttt{9\_12\_4\_multiscale/epoch\_15.pth} and CLIP ViT-L/14@336px. Images are resized and center-cropped to $518\times518$ with OpenAI CLIP normalization. Cached features are extracted from layers $\{6,12,18,24\}$, while controlled WPTA ablations aggregate feature map layers $\{1,2,3\}$. Layer fusion uses summation, Gaussian smoothing uses $\sigma=5$, and AUPRO uses 200 thresholds.

Prototype adaptation uses temperature $0.07$, semantic sharpening $\gamma=1.0$, wavelet evidence exponent $\eta=1.0$, and evidence top-$k$ ratio $0.2$. Boundary-aware wavelet reliability uses boundary-aware mode with mix weight $0.05$. Except for the fixed-baseline row, prototype/fusion variants in controlled ablations share the same multi-crop and pixel-to-image settings.

\subsection{Results}

\subsubsection{Industrial Results}
Table~\ref{tab:industry_result} follows the compact multi-method layout of the referenced LNCS paper. It reports representative CLIP-based ZSAD methods together with our current fixed AnomalyCLIP baseline and WPTA results. External rows are used as protocol-reference positioning; the same-protocol claim remains the improvement over the current fixed AnomalyCLIP baseline and the controlled ablations below.

\begin{table}[htbp]
  \centering
  \caption{Zero-shot anomaly detection results on five industrial benchmarks. Each cell reports paired metrics. Best and second-best values among collected numbers are marked in \textbf{bold} and \underline{underline}. $^{\dagger}$ denotes reproduced results in the referenced table, and $^{\ddagger}$ denotes our current fixed-baseline evaluation.}
  \label{tab:industry_result}
  \resizebox{\textwidth}{!}{%
    \begin{tabular}{llllllllll}
      \toprule
      \multicolumn{10}{c}{\textbf{Image-level results (Image-AUROC, Image-AP)}} \\
      \midrule
      Dataset & CLIP & WinCLIP & VAND & CoOp & AdaCLIP & AnomalyCLIP$^{\ddagger}$ & AA-CLIP$^{\dagger}$ & TAAP/INPC & WPTA \\
      \midrule
      MVTec AD & (74.1, 87.6) & (91.8, 96.5) & (86.1, 93.5) & (88.8, 94.8) & (89.2, 96.4) & (91.6, 96.4) & (90.5, 94.9) & (\underline{92.8}, \underline{96.7}) & (\textbf{94.5}, \textbf{97.6}) \\
      VisA & (66.4, 71.5) & (78.1, 81.2) & (78.0, 81.4) & (62.8, 68.1) & (\textbf{85.8}, 84.9) & (82.0, 85.3) & (\underline{84.6}, 82.2) & (83.3, \underline{85.9}) & (\underline{84.6}, \textbf{87.4}) \\
      MPDD & (54.3, 65.4) & (63.6, 69.9) & (73.0, 80.2) & (55.1, 64.2) & (76.0, 80.4) & (73.7, 76.5) & (75.1, 80.1) & (\textbf{81.8}, \textbf{84.8}) & (\underline{77.8}, \underline{82.3}) \\
      BTAD & (34.5, 52.5) & (68.2, 70.9) & (73.6, 68.6) & (66.8, 77.4) & (88.6, 92.4) & (89.1, 91.0) & (\textbf{94.8}, \textbf{97.9}) & (93.3, 93.1) & (\underline{93.9}, \underline{94.9}) \\
      DTD-Synthetic & (71.6, 85.7) & (93.2, 92.6) & (86.4, 95.0) & (-, -) & (95.5, 97.0) & (94.5, 97.7) & (93.3, 97.8) & (\textbf{97.5}, \textbf{99.0}) & (\underline{96.9}, \underline{98.7}) \\
      \midrule
      \textbf{Average} & (60.2, 72.5) & (79.0, 82.2) & (79.4, 83.7) & (68.4, 76.1) & (87.0, 90.2) & (86.2, 89.4) & (87.7, 90.6) & (\textbf{89.7}, \underline{91.9}) & (\underline{89.5}, \textbf{92.2}) \\
      \midrule
      \multicolumn{10}{c}{\textbf{Pixel-level results (Pixel-AUROC, Pixel-AUPRO)}} \\
      \midrule
      MVTec AD & (38.4, 11.3) & (85.1, 64.6) & (87.6, 44.0) & (33.3, 6.7) & (88.7, 37.8) & (91.2, 83.2) & (\textbf{91.9}, 84.6) & (91.5, \underline{85.5}) & (\underline{91.8}, \textbf{85.6}) \\
      VisA & (46.6, 14.8) & (79.6, 56.8) & (94.2, 86.8) & (24.2, 3.8) & (95.5, 56.8) & (95.5, 86.7) & (95.5, 83.0) & (\underline{95.6}, \underline{88.3}) & (\textbf{96.2}, \textbf{91.3}) \\
      MPDD & (62.1, 33.0) & (76.4, 48.9) & (94.1, 83.2) & (15.4, 2.3) & (96.1, 60.3) & (\underline{96.9}, 84.6) & (96.7, 76.5) & (96.7, \textbf{89.9}) & (\textbf{97.3}, \textbf{89.9}) \\
      BTAD & (30.6, 4.4) & (72.7, 27.3) & (60.8, 25.0) & (28.6, 3.8) & (92.1, 32.5) & (93.5, 70.5) & (\textbf{97.0}, 69.0) & (94.1, \underline{77.1}) & (\underline{96.3}, \textbf{78.2}) \\
      DTD-Synthetic & (33.9, 12.5) & (83.9, 57.8) & (95.3, 86.9) & (-, -) & (97.7, 75.0) & (97.4, 89.1) & (96.4, 85.9) & (97.4, \underline{91.5}) & (\textbf{97.9}, \textbf{91.8}) \\
      \midrule
      \textbf{Average} & (42.3, 15.2) & (79.5, 51.1) & (86.4, 65.2) & (25.4, 4.2) & (94.0, 52.5) & (94.9, 82.8) & (\underline{95.5}, 79.8) & (95.1, \underline{86.5}) & (\textbf{95.9}, \textbf{87.4}) \\
      \bottomrule
    \end{tabular}%
  }
\end{table}

\subsubsection{Core Component Ablation}
Table~\ref{tab:core_ablation} validates the WPTA mechanism on MVTec and VisA. Direct wavelet fusion is worse than the fixed baseline, showing that wavelet responses should not be used as final anomaly scores. Semantic prototype adaptation improves over fixed prototypes, while wavelet-guided adaptation and conservative update further improve Pixel-AUPRO.

\begin{table}[htbp]
  \centering
  \caption{Core component ablation on MVTec and VisA. Each cell reports (P-AUROC, P-AUPRO) / (I-AUROC, I-AP).}
  \label{tab:core_ablation}
  \resizebox{\textwidth}{!}{%
    \begin{tabular}{lll}
      \toprule
      Method & MVTec AD & VisA \\
      \midrule
      Baseline & (91.2, 83.2) / (91.6, 96.4) & (95.5, 86.7) / (82.0, 85.3) \\
      Direct wavelet fusion / no adaptation & (88.7, 80.4) / (92.9, 96.9) & (94.6, 85.1) / (81.6, 84.8) \\
      Semantic prototype adaptation & (91.6, 85.2) / (93.7, 97.1) & (96.0, 90.4) / (83.7, 86.9) \\
      Wavelet prototype adaptation w/o conservative & (91.7, 85.8) / (93.9, 97.2) & (96.1, 91.3) / (84.1, 87.0) \\
      Full WPTA controlled setting & \textbf{(91.8, 86.2) / (94.1, 97.4)} & \textbf{(96.2, 91.7) / (84.3, 87.3)} \\
      \bottomrule
    \end{tabular}%
  }
\end{table}

\subsubsection{Wavelet Design Ablation}
Table~\ref{tab:wavelet_design} studies the role of wavelet reliability. HF-only reliability gives limited gains, while boundary-aware reliability improves Pixel-AUPRO by suppressing structure-boundary pseudo evidence.

\begin{table}[htbp]
  \centering
  \caption{Wavelet reliability design ablation. Each cell reports (P-AUROC, P-AUPRO) / (I-AUROC, I-AP).}
  \label{tab:wavelet_design}
  \resizebox{\textwidth}{!}{%
    \begin{tabular}{lll}
      \toprule
      Wavelet setting & MVTec AD & VisA \\
      \midrule
      Semantic-only prototype adaptation & (91.6, 85.2) / (93.7, 97.1) & (96.0, 90.4) / (83.7, 86.9) \\
      Direct wavelet fusion & (88.7, 80.4) / (92.9, 96.9) & (94.6, 85.1) / (81.6, 84.8) \\
      HF-only W + prototype adaptation & (91.6, 85.3) / (94.0, 97.2) & (96.0, 90.8) / (84.0, 86.9) \\
      Boundary-aware W + prototype adaptation & (91.7, 85.7) / (93.8, 97.3) & (96.1, 91.2) / (83.9, 87.1) \\
      Full boundary-aware W + conservative & \textbf{(91.8, 86.2) / (94.1, 97.4)} & \textbf{(96.2, 91.7) / (84.3, 87.3)} \\
      \bottomrule
    \end{tabular}%
  }
\end{table}

\subsubsection{Configuration Audit}
\mbox{}\\
\noindent
Table~\ref{tab:final_config} records which modules are active in each dataset-specific final system. It prevents over-attributing all five-dataset gains to the WPTA mechanism.

\begin{table}[htbp]
  \centering
  \caption{Configuration audit for the five-dataset calibrated inference system.}
  \label{tab:final_config}
  \resizebox{\textwidth}{!}{%
    \begin{tabular}{lccccc}
      \toprule
      Dataset & Result & Wavelet reliability & Prototype rectification & Multi-crop & Pixel-to-image \\
      \midrule
      MVTec AD & 91.8 / 85.6 / 94.5 / 97.6 & yes & yes & yes & yes \\
      VisA & 96.2 / 91.3 / 84.6 / 87.4 & yes & yes & yes & yes \\
      MPDD & 97.3 / 89.9 / 77.8 / 82.3 & no & no & yes & yes \\
      BTAD & 96.3 / 78.2 / 93.9 / 94.9 & no & no & yes & yes \\
      DTD-Synthetic & 97.9 / 91.8 / 96.9 / 98.7 & yes & no & yes & yes \\
      \bottomrule
    \end{tabular}%
  }
\end{table}

\section{Discussion}

\paragraph{Why wavelet reliability, not wavelet scoring.}\mbox{}\\
\noindent
Wavelets provide local frequency structure rather than semantic anomaly decisions. The direct-fusion negative control shows that adding wavelet response to the anomaly map hurts localization. WPTA instead uses wavelets only as evidence reliability and lets calibrated CLIP prototypes produce the final anomaly map.

\paragraph{Why conservative calibration.}\mbox{}\\
\noindent
Unlabeled test-time calibration can suffer from prototype drift. WPTA uses confidence-gated conservative updates so that visual anchors affect prototypes only when the evidence is reliable. In the default setting, the abnormal prototype update strength is zero, and abnormal evidence mainly controls the update gate. This makes the method conservative rather than a large bidirectional prototype shift.

\paragraph{Scope and limitations.}\mbox{}\\
\noindent
The five-dataset final systems use dataset-specific configurations and therefore support system-level gains over the fixed AnomalyCLIP baseline, not a causal claim that WPTA is isolated on all five datasets. The WPTA mechanism is validated by controlled MVTec/VisA ablations. External method comparisons require careful protocol alignment before being used for ranking claims.

\section{Conclusion}

We proposed WPTA for CLIP-based zero-shot anomaly detection. The central idea is to use Haar wavelet cues as patch evidence reliability for test-time prototype adaptation instead of direct final-map fusion. Controlled MVTec/VisA ablations show that direct wavelet fusion hurts localization, while semantic-spectral evidence selection and conservative prototype calibration improve test-time adaptation. Across five industrial benchmarks, the calibrated inference system improves the fixed AnomalyCLIP baseline, especially on Pixel-AUPRO. These results support a conservative conclusion: wavelet cues are more useful as evidence reliability than as anomaly scores, and test-time prototype calibration can reduce the mismatch between fixed CLIP prototypes and instance-specific defect appearances.

\bibliographystyle{splncs04}
\bibliography{wpta_references_v0.1}

\end{document}
